\section{Method}
\label{sec:method}

\paragraph{Task formulation.}
Each Vietnamese social-media comment $x_i$ is associated with three
groups of labels. The \emph{pragmatic} group contains six binary
fields $\mathbf{p}_i\in\{0,1\}^6$: implicit sentiment, sarcasm, irony,
idiom/figurative, code-switching, and mocking. The \emph{polarity}
group contains a 3-way signed sentiment label
$s_i\in\{\text{negative}, \text{neutral}, \text{positive}\}$ over the
\emph{intended} meaning (so a sarcastic ``\textit{hay l\d{a}m}'' is
labelled negative). The \emph{emotion} group is the
UIT-VSMEC-compatible 7-way emotion label $e_i$. We treat all three
groups as a single multi-task target.

\paragraph{Backbones.}
\method{} is backbone-agnostic. We use three backbones in the main
experiments. (i) \textbf{PhoBERT-base}
\citep{phobert2020}, the default encoder for Vietnamese, fine-tuned
end-to-end. (ii) \textbf{XLM-R-large}
\citep{conneau2020xlmr}, used to measure multilingual transfer.
(iii) \textbf{Vistral-7B} \citep{vistral2024}, a small Vietnamese
instruction LM, with LoRA \citep{hu2022lora} on attention projections.
For Sailor-7B \citep{sailor2024} we use the same LoRA recipe. We do not
treat the Vistral-7B variant as the default because its annotation
cost is the same as PhoBERT's but its training cost is $\approx 5\times$
higher (Table~\ref{tab:cost}).

\paragraph{Multi-task heads.}
Let $h_i\in\mathbb{R}^d$ be the comment representation from the
backbone (the \texttt{[CLS]} embedding for PhoBERT and XLM-R, the
mean pooled last-layer representation for Vistral/Sailor). \method{}
attaches:
\begin{itemize}\itemsep1pt
  \item Six pragmatic binary heads $g^{\text{prag}}_k(h_i)$ trained
    with binary cross-entropy on $\mathbf{p}_i^{(k)}$.
  \item One signed-polarity head $g^{\text{pol}}(h_i)$ trained with
    cross-entropy on $s_i$.
  \item One emotion head $g^{\text{emo}}(h_i)$ trained with
    cross-entropy on $e_i$.
  \item One explanation head $g^{\text{exp}}$, a small Transformer
    decoder of $2$ layers and $128$ hidden units that consumes $h_i$
    and is trained with token-level cross-entropy on a 1--2 sentence
    Vietnamese rationale $r_i$ produced once from GPT-4o-mini with
    the gold pragmatic labels included in the prompt. The decoder is
    used only during training and is dropped at inference.
\end{itemize}

\paragraph{Loss with uncertainty weighting.}
We combine the heads with the homoscedastic-uncertainty weighting of
\citet{kendall2018uncertainty}: for each task $t$ with loss
$\mathcal{L}_t$ we introduce a learnable log-variance $\sigma_t$ and
optimise
\[
\mathcal{L} \;=\; \sum_t \frac{1}{2\sigma_t^2}\mathcal{L}_t
              \;+\; \tfrac{1}{2}\log\sigma_t^2 \;+\;
              \beta\,\mathcal{L}^{\text{exp}},
\]
where $\mathcal{L}^{\text{exp}}$ is the token-level cross-entropy of
the rationale and $\beta{=}0.3$ is a small fixed weight. This avoids
hand-tuning per-task weights while still letting the rare pragmatic
tasks dominate.

\paragraph{Explanation-augmented training.}
For every gold-labelled training comment we collect a short
Vietnamese rationale from GPT-4o-mini under a system prompt that
states (i) the comment text, (ii) the gold pragmatic and polarity
labels, and (iii) the instruction to produce a 1--2 sentence
explanation in Vietnamese that names the lexical or contextual cues
responsible for the labels. Annotators then review a 5\% sample for
faithfulness; rationales judged unfaithful are dropped. The
remaining rationales become the auxiliary target for $g^{\text{exp}}$,
following the rationale-distillation recipe of
\citet{hsieh2023distillstep}, \citet{magister2023teaching} and
\citet{shridhar2023distillingcot}. At inference, the rationale head
is removed, so prediction is a single forward pass through the
backbone and classification heads.

\paragraph{Direct classification vs.\ explanation-then-classification.}
We compare three regimes. \textbf{Direct}: heads predict labels
directly from $h_i$ and the rationale head is removed at inference.
\textbf{CoT-only}: the rationale head replaces the classification
heads at inference; labels are read off the rationale text with a
deterministic post-processor. \textbf{ViPragSent (default)}: both
training signals are used, but inference uses the classification
heads. \method{} dominates both single regimes
(Table~\ref{tab:main_pragmatic}); the value of the rationale head is
that it shapes the encoder, not that it generates a runtime
explanation.

\paragraph{Inference-time cost.}
Because the rationale head is dropped at inference, \method{}'s
runtime cost is exactly that of the underlying backbone with seven
small classification heads. On PhoBERT-base, end-to-end latency is
$5.1$ ms per comment on a single A100 (batch $64$); on Vistral-7B it
is $34.7$ ms per comment with LoRA. We report compute and annotation
costs in Table~\ref{tab:cost}.

\paragraph{Artifact contract.}
Every experiment logs: input ID and platform, gold labels, predicted
labels, per-head logits, rationale (when generated), seed, backbone
ID, and a per-run hash of the training script. The result JSON files
in \texttt{results/} can be regenerated end-to-end from
\texttt{scripts/make\_artifacts.py}, and the claim ledger
\texttt{results/claim\_ledger.csv} links each table and figure
anchor to its result file.
